\section{Problem Formulation and Methodology}
\label{sec:problem_method}
\label{sec:problem_formulation}

Let $\mathcal{C}$ be the set of cells and let $t$ index hourly timestamps. For each cell-time pair $(c,t)$, the panel contains a KPI/business vector $\mathbf{x}_{c,t}\in\mathbb{R}^{M}$, an alarm/fault vector $\mathbf{a}_{c,t}\in\mathbb{R}^{A}$, an anonymized cell identifier and timestamp-derived features. In the evaluated anonymized MBB operations panel, $|\mathcal{C}|=33$, $M=48$ and $A=137$, although only 11 alarm/fault columns are active in the observed slice. Given a lookback window $\mathcal{W}_{c,t}=\{(\mathbf{x}_{c,\tau},\mathbf{a}_{c,\tau}):t-L+1\leq \tau\leq t\}$, the objective is to predict outcomes at horizon $h\in\{1,3,6,12\}$ hours. Windows crossing the detected long time gap are removed before modeling.

A feature representation for target $(c,t+h)$ is called temporally admissible if it is constructed only from $\mathcal{W}_{c,t}$, static or timestamp-derived identifiers available at $t$, and statistics fitted on the chronological training partition. CET-STGL uses this constraint throughout to separate operationally available evidence from future outcomes.

\subsection{KPI Forecasting}
The forecasting task is to predict future KPI values, as formalized here below:
\begin{equation}
    \hat{\mathbf{x}}_{c,t+h} = f_h(\mathcal{W}_{c,t}, c, t),
\end{equation}
where $\hat{\mathbf{x}}_{c,t+h}$ is the predicted KPI vector at horizon $h$, $\mathcal{W}_{c,t}$ is the historical lookback window for cell $c$ at time $t$, and $f_h$ represents the forecasting mapping function. Forecasting labels are real future observations from the private cellular panel. We report aggregate metrics over all KPI/business columns and paper-level metrics over the primary KPI subset.

\subsection{Weak Degradation Classification}
The degradation task is to predict a weak binary label, formalized as follows:
\begin{equation}
    \hat{y}_{c,t+h} = g_h(\mathcal{W}_{c,t}, c, t), \quad y_{c,t+h}\in\{0,1\},
\end{equation}
where $\hat{y}_{c,t+h}$ is the predicted degradation probability, $y_{c,t+h}$ is the weak binary indicator, and $g_h$ denotes the classification function. The label is derived from thresholds fitted on the training split. The evaluated implementation uses five independent OR conditions: high call-drop rate above the larger of 1.0\% and the cell-specific training 99th percentile; low Radio Resource Control (RRC) setup success rate below the smaller of 98.0\% and the cell-specific training 5th percentile; low Evolved Radio Access Bearer (E-RAB) setup success rate below the smaller of 98.0\% and the cell-specific training 5th percentile; low downlink user throughput in Mbps below the cell-specific training 10th percentile; or low dimensionless Channel Quality Indicator (CQI) below the cell-specific training 10th percentile. Cell-specific thresholds are replaced by the corresponding full-training-partition thresholds when fewer than 24 non-missing training samples are available for the relevant cell and KPI, or when the computed cell-level threshold is not finite. The label is therefore a weak KPI degradation proxy rather than an expert incident label, following weak-supervision practice for indirect or noisy supervision \cite{zhou2018weakly,ratner2020snorkel}.

\subsection{Weak Alarm-Candidate Prioritization}
For degraded samples with recent alarms, the ranking task scores candidate alarm types $a\in\mathcal{A}_{c,t}$ as follows:
\begin{equation}
    s_h(a,c,t)=r_h(a,\mathcal{W}_{c,t},c,t),
\end{equation}
where $s_h(a,c,t)$ is the prioritization score assigned to alarm candidate $a$ in cell $c$ at time $t$, and $r_h$ is the scoring function. The target candidate is induced by train-only alarm-to-degradation lift and recent event intensity. The ranking task evaluates weak-label consistency only. It is not expert-confirmed root-cause localization.

\subsection{Causal-Inspired Graph Boundary}
The event influence graph is constructed from temporal precedence and train-only association. It is useful for organizing candidate event influence, but the observational data contain no interventions or expert causal annotations. We therefore use only causal-inspired terminology and avoid causal-effect claims.
